% YAAC Another Awesome CV LaTeX Template
%
% This template has been downloaded from:
% https://github.com/darwiin/yaac-another-awesome-cv
%
% Author:
% Christophe Roger
%
% Template license:
% CC BY-SA 4.0 (https://creativecommons.org/licenses/by-sa/4.0/)
%Section: Work Experience at the top
\sectionTitle{Experiencia Profesional}{\faSuitcase}
\begin{experiences}
% SEIDOR se escribe como filas sueltas de longtable (en vez de usar
% \experience, que mete el cuerpo entero en un minipage no divisible) porque
% en español el bloque de 5 clientes es demasiado alto para caber en el
% espacio que sobra en la pagina 1, y longtable nunca parte una unica fila
% entre paginas: con \experience, todo el bloque saltaria entero a la
% pagina 2 dejando un hueco enorme en la 1. Como filas independientes,
% longtable reparte los \clientblock entre paginas de forma natural.
\textbf{Junio 2025} & \textbf{Ingeniero de Machine Learning Senior - Consultor, \textsc{SEIDOR ANALYTICS}, Lima, PE} \hfill \emph{Híbrido, Completo} \\[\cvrowheadsep]
\textbf{Presente} & Brindé servicios de consultoría para empresas del sector bebidas en toda LATAM, diseñando plataformas de datos escalables y sistemas de ML end-to-end. \\
    & \clientblock{Koandina}{Industria de Bebidas}{%
      \begin{itemize}
        \item Construí arquitecturas de datos en tiempo real usando Kinesis, Glue y Athena, exponiendo datasets curados mediante APIs de AppSync.
        \item Entregué soluciones de API y streaming para equipos de ML y analítica en operaciones multi-país.
        \item Implementé Route 53 para escalamiento de DNS y recuperación ante desastres.
        \item Optimicé la infraestructura cloud mediante la refactorización de módulos de ingesta de datos legados y automatización de despliegues con IaC (CDK/SAM), reduciendo costos operativos.
      \end{itemize}
    } \\
    & \clientblock{Bodegas La Negrita (BLN)}{Distribución de Licores}{%
      \begin{itemize}
        \item Lideré el desarrollo de modelos de pronóstico en Snowflake ML para múltiples líneas de producto.
        \item Diseñé pipelines de features y flujos de inferencia automatizados.
        \item Construí dashboards en Streamlit y Power BI; integré datos de Azure en Snowflake para analítica e ingeniería de features.
      \end{itemize}
    } \\
    & \clientblock{Arca Continental}{Industria de Bebidas}{%
      \begin{itemize}
        \item Implementé un framework MLOps end-to-end en Snowflake para modelos de regresión.
        \item Usé Feature Store y Model Registry para la gobernanza del ciclo de vida.
        \item Aproveché el escalamiento de clústeres para paralelizar entrenamiento de modelos a gran escala.
      \end{itemize}
    } \\
    & \clientblock{Industrias San Miguel (ISM)}{Industria de Bebidas}{%
      \begin{itemize}
        \item Desarrollé soluciones de pronóstico de demanda en Snowflake ML.
        \item Diseñé features avanzadas en Python y apliqué clustering para mejorar la precisión de los pronósticos.
      \end{itemize}
    } \\
    & \clientblock{Innova Sports}{Industria Retail}{%
      \begin{itemize}
        \item Desarrollé soluciones de pronóstico de demanda para productos retail basadas en escenarios de campañas de descuento planificadas.
        \item Usé notebooks de Amazon SageMaker para análisis exploratorio de datos, training jobs de SageMaker para el desarrollo de modelos, y SageMaker Pipelines para automatizar el entrenamiento e inferencia de modelos.
      \end{itemize}
    } \\
    & \rule{0pt}{\cvtagrowheight}\small{\foreach \n in {Snowflake ML, {AWS (Kinesis, Athena, SageMaker, CloudFormation)}, Azure, Python, MLOps}{\cvtag{\n}}}\vspace{\cvrowtagsep} \\

\experience
  {Noviembre 2024}{Tech Lead}{Sokso}{Lima, PE}
  {Híbrido}{Completo}
  {Junio 2025} {
    \begin{itemize}
      \item Liderando el diseño arquitectónico de la plataforma SMART 3.0, supervisando y asignando tareas a los equipos de desarrollo backend y frontend, y estableciendo estándares técnicos de implementación.
      \item Encabezando el proyecto estratégico de migración del sistema ERP de SAP a NetSuite, coordinando la transferencia de datos críticos, rediseñando procesos de negocio y asegurando la continuidad operativa durante la transición.
      \item Gestionando la infraestructura AWS de la empresa, implementando automatizaciones de despliegue en instancias EC2, configurando pipelines de CI/CD y optimizando costos y rendimiento de los servicios cloud.
    \end{itemize}
  }
  {{AWS (EC2, ECS, S3)}, CI/CD, DevOps, NetSuite, SAP, ERP, Nest.js, Next.js, Technical Leadership}

\experience
  {Noviembre 2023}{Desarrollador Full Stack}{Sokso}{Lima, PE}
  {Híbrido}{Completo}
  {Noviembre 2024} {
    \begin{itemize}
      \item Lideré el diseño e implementación de mejoras avanzadas al frontend (Vue.js) de la plataforma SMART 2.0, optimizando la experiencia de usuario en módulos críticos como el sistema de contacto con clientes, y apliqué mejoras en prácticas de ciberseguridad, estableciendo estándares de desarrollo seguro tanto para frontend como backend (.NET).
      \item Contribuí al desarrollo inicial de la versión 3.0 de la plataforma core SMART, usando Nest.js en el backend, Next.js en el frontend, y configurando un pipeline DevOps robusto en AWS (ECS, S3, CI/CD), logrando una infraestructura ágil, escalable y de alta disponibilidad.
      \item Implementé múltiples automatizaciones en Python para agilizar tareas en las áreas comercial y de diseño, incluyendo scripts de procesamiento de datos, generación de reportes personalizados y flujos de trabajo automatizados que mejoraron la eficiencia y redujeron el tiempo de trabajo manual en ambas áreas.
      \item Desarrollé proyectos piloto de machine learning para pronóstico de ventas de los próximos meses y un modelo predictivo de churn de clientes usando técnicas avanzadas de series de tiempo (Python, PyTorch, GCP), lo que mejoró la precisión en la gestión de inventario y permitió anticipar patrones de abandono, optimizando la retención de clientes.
    \end{itemize}
  }
  {Vue.js, .NET, SQL Server, Machine Learning, Python, GCP, {AWS (EC2, ECS, S3)}, CI/CD, Next.js, DevOps}

\emptySeparator
\experience
  {Noviembre 2023}{Desarrollador Full Stack especializado en IA}{Arkrisk}{New York, US}
  {Remoto}{Parcial}
  {Marzo 2025} {
    \begin{itemize}
      \item Desarrollé e implementé un sistema RAG (Retrieval-Augmented Generation) con memoria, function tools e historial de usuario usando FastAPI, LangGraph y Ollama, integrando modelos como DeepSeek, Llama3 y Phi4 para garantizar la privacidad de los datos. El sistema fue desplegado en una instancia de AWS con GPU.
      \item Implementé múltiples soluciones LLM en producción: usando Hugging Face Inference Endpoints para modelos pre-entrenados y desarrollando modelos personalizados con PyTorch desplegados en instancias EC2 de AWS. Integré estos modelos con tecnología OCR para la digitalización y análisis de pólizas de seguro, permitiendo la detección de patrones de comportamiento en los datos y optimizando la eficiencia operativa.
      \item Implementé automatizaciones en Python para optimizar procesos internos, creando flujos de trabajo que agilizaron el procesamiento de datos y la generación de reportes.
      \item Realicé integraciones con servicios cloud de AWS, usando DynamoDB y S3 para almacenar y gestionar grandes volúmenes de datos de seguros y clientes.
      \item Desarrollé y optimicé el frontend con Next.js, creando vistas que consumen APIs de FastAPI, mejorando la experiencia de usuario y la interacción con los modelos desplegados.
      \item Trabajé en integraciones con \textbf{OpenAI assistant v2} y desarrollé soluciones basadas en RAG para chatbots de QA, mejorando la capacidad de respuesta y la calidad del servicio al cliente.
    \end{itemize}
  }
  {FastAPI, LangGraph, Ollama, DeepSeek, Hugging Face, OCR, Python, {AWS (EC2, DynamoDB)}, Next.js, LLMs, OpenAI, RAGs, Deep Learning}

\emptySeparator
\experience
  {Mayo 2023}{Ingeniero de Machine Learning}{Impacto AI}{Lima, PE}
  {Remoto}{Completo}
  {Octubre 2023} {
    \begin{itemize}
      \item Lideré el desarrollo de Chatbots para WhatsApp, Messenger e Instagram basados en modelos avanzados como Chat-GPT, mejorando la comprensión y las respuestas naturales. También usé OCR basado en OpenCV avanzado, modelos de generación de imágenes y modelos de reconocimiento de voz.
      \item Lideré proyectos de Data Science, enfocados en análisis de datos y desarrollo de modelos ML usando Python y tecnologías Big Data como PySpark sobre AWS.
      \item Automaticé despliegues de modelos ML en producción usando Airflow y gestioné pipelines de datos en tiempo real y por lotes.
    \end{itemize}
  }
  {Chat-GPT, Spark, SQL, MLOps, AWS, Big Data, OpenCV}

\emptySeparator
\experience
  {Enero 2022}{Desarrollador Full Stack}{AmauTeach Project}{Lima, PE}
  {Remoto}{Completo}
  {Diciembre 2022} {
    \begin{itemize}
      \item Desarrollé AmauTeach Exam, una plataforma de exámenes en línea enfocada en la prevención de plagio usando técnicas avanzadas de procesamiento de datos y machine learning. Usé tecnologías como Vue, Flask, FastAPI y Node.
      \item Participé en el diseño de sistemas de visualización de datos para comunicar insights de forma efectiva a los stakeholders.
      \item Realicé integraciones con servicios de AWS, usando Cognito para autenticación de usuarios, S3 para almacenamiento de contenido multimedia, y funciones Lambda para optimizar microservicios relacionados con el procesamiento de datos.
    \end{itemize}
  }
  {Vue, Node.js, Security, AWS, Big Data}

\emptySeparator
\experience
  {Julio 2021}{Desarrollador Full Stack}{Chazki}{Lima, PE}
  {Remoto}{Completo}
  {Diciembre 2021} {
    \begin{itemize}
      \item Contribuí al desarrollo de funcionalidades críticas en la plataforma logística de Chazki, trabajando en servicios web y la aplicación móvil, usando tecnologías como Vue, React Native, Node y Google Cloud para visualizar y aprovechar insights generados a partir de los datos.
      \item Contribuí al backend desarrollado con Spring Boot para el cálculo de rutas mínimas de entrega para los conductores, implementando algoritmos de optimización con Java, usando Spring Data JPA y REST APIs, logrando mejores tiempos de entrega y costos operativos optimizados.
      \item Implementé y mantuve brokers de mensajería como Kafka para asegurar una transmisión de datos eficiente y en tiempo real en un entorno Big Data.
    \end{itemize}
  }
  {Logistics, Vue, React Native, Node, Java, Spring Boot, Spring Data JPA, REST APIs, Databases, Big Data, Kafka}

\emptySeparator
\experience
  {Octubre 2020}{Desarrollador Full Stack y Científico de Datos}{MDP CONSULTING}{Lima, PE}
  {Remoto}{Completo}
  {Julio 2021} {
    \begin{itemize}
      \item Gestioné pipelines de datos y servicios de IA en IBM Cloud para mejorar la gestión de cobranza en instituciones bancarias, asegurando la seguridad de los datos en entornos Big Data.
      \item Desarrollé modelos de Machine Learning para brindar insights sobre el uso del servicio y su impacto en los usuarios finales, usando librerías como pandas, scikit-learn y Jupyter Notebooks.
      \item Realicé pruebas exhaustivas de seguridad, incluyendo pruebas de inyección SQL con herramientas como SQLInjectMe, Sqliv y Sqlmap, como medida preventiva ante posibles ataques.
    \end{itemize}
  }
  {Data Security, Flask, Automation, Big Data}
\end{experiences}
